\documentclass[epj]{svjour}

\usepackage[T1]{fontenc}
\usepackage{lmodern}
\usepackage{amsmath,amssymb,amsfonts}
\usepackage{mathtools}
\usepackage{bm}
\usepackage{graphicx}
\usepackage{booktabs}
\usepackage{multirow}
\usepackage{array}
%\usepackage{xurl}      
\usepackage{hyperref}
\usepackage{enumitem}

\graphicspath{{figures/}}

\tolerance=2000
\emergencystretch=3em
\raggedbottom

\begin{document}

% PTUM_LIBRARY_FINAL_v1.1_CLEAN
% Self-contained library-final manuscript source.
% This version removes the obsolete external-supplement pointer,
% moves implementation details into the appendix, and preserves
% explicit claim boundaries for later submission work.

\title{Phase Transition Universes Model:\\
A Discharge-Based Cosmological Creation Scenario\\
with a Unified Dark Sector}

\titlerunning{PTUM: A Discharge-Based Cosmological Creation Scenario}

\authorrunning{E. Esa}

\author{Erdem Esa}

\offprints{E. Esa}

\institute{Independent Researcher, Istanbul, Turkey\\
\email{erdem.esa.71@gmail.com}}

\date{}

\abstract{%
The dark sector of the standard cosmological model comprises dark matter, dark energy, and possible dark radiation, whose common origin remains unexplained. We present the Phase Transition Universes Model (PTUM), a discharge-based cosmological creation scenario where a scalar quantum fuel field $\Phi$ accumulates, reaches a critical threshold, and undergoes a tachyonic transition. Within this effective description, the discharge event leaves behind a dark-matter candidate sector ($\chi$), a residual dark-energy sector ($\Phi_{\rm res}$), and a relativistic proxy-fluid dark radiation component (pfd, $\psi$). A falsifiable prediction of PTUM is that, within the effective linear proxy-fluid description, the pfd component carries $w=1/3$, $c_s^2=1/3$, and zero anisotropic stress ($\sigma=0$). This produces no neutrino-induced CMB acoustic peak phase shift, distinguishing it from free-streaming dark radiation---a signature that may be tested at high significance by next-generation CMB experiments such as CMB-S4, depending on the realized pfd abundance. We implement the effective phenomenological system in \texttt{CLASS} at background and perturbation levels. The creation sector supports a plateau-type slow-roll phase compatible with current CMB constraints, fixing the discharge scale to $V_0^{1/4}\sim 10^{15}$--$10^{16}\,\mathrm{GeV}$. The dark-energy sector admits a broad near-$\Lambda$ viable region with $w_{\rm DE}\approx -1$. A Bayesian benchmark analysis using Planck 2018, BAO, and Pantheon+ data yields $\Delta N_{\rm fluid} < 0.31$ at $95\%$ credibility. Model comparison gives $\Delta\chi^2 = -0.26$ and $\Delta\mathrm{AIC}=2.26$; current data therefore show a weak preference for $\Lambda$CDM while leaving PTUM viable as a falsifiable effective extension. Within the scanned benchmark regime, the sectors satisfy EFT validity, BBN consistency, and linear stability. The formulation should be read as an effective phenomenological framework rather than as a first-principles derivation of the pre-discharge background, a solution of the cosmological-constant hierarchy, or a UV completion.%
} % abstract'ın kapanış parantezi
\PACS{
      {98.80.Cq}{Particle-theory and field-theory models of the early Universe} \and
      {95.36.+x}{Dark energy} \and
      {95.35.+d}{Dark matter} \and
      {98.80.Es}{Observational cosmology}
     } % end of PACS
\maketitle

%============================================================
\section{Introduction}
\label{sec:intro}
%============================================================

The standard cosmological model accounts for the later evolution of
the Universe through a combination of radiation, baryons, cold dark
matter, and a cosmological constant~\cite{Planck2018params}, yet it
does not explain why the Universe began in a hot, dense state. More
pointedly, the dark sector---which dominates the cosmic energy
budget---is composed of at least two phenomenologically distinct
components whose microscopic origins are unrelated within
$\Lambda$CDM. Dark matter is introduced as a pressureless fluid
needed to explain structure formation and galactic rotation curves;
dark energy is introduced as a negative-pressure component needed to
explain the accelerating expansion. Whether these components share a
common physical origin, and whether the hot Big Bang itself can be
understood as a dynamical event rather than an initial condition, are
among the deepest open questions in theoretical cosmology.

Several approaches to these questions have been explored. Quintessence
and dynamical dark-energy models replace the cosmological constant
with a slowly rolling scalar field~\cite{RatraPeebles,Caldwell1998},
while unified dark-sector models attempt to describe dark matter and
dark energy through a single fluid~\cite{Kamenshchik2001,Bento2002}.
Extra relativistic species parametrised by
$\Delta N_{\rm eff}$~\cite{Mangano2005,Cyburt2016} probe the
radiation content of the early Universe. However, none of these
frameworks addresses the origin of the hot expanding state itself, nor
do they predict specific relationships between the dark-sector
components. The question of \emph{why} the dark sector has its
observed composition remains unanswered.

The Phase Transition Universes Model (PTUM) was developed to address
these issues from a different starting point. Instead of treating dark
matter and dark energy as independent sectors added to a pre-existing
cosmological framework, PTUM proposes a discharge-based creation
mechanism in which \emph{all} dark-sector components arise from a
single cosmogonic event. In this picture, a scalar quantum fuel field
$\Phi$ accumulates in a metastable multiverse-like background and,
once a critical threshold is reached, undergoes a tachyonic transition
that activates string-like proxy modes. The resulting discharge is modeled as
separating into:

\begin{itemize}
\item a clustering passive component $\rightarrow$ a dark-matter candidate sector~($\chi$),
\item a residual negative-pressure component $\rightarrow$
      dark energy~($\Phi_{\rm res}$),
\item a radiation-like transient component $\rightarrow$
      proxy-fluid dark radiation~(pfd, $\psi$).
\end{itemize}

The most distinctive---and directly falsifiable---observable
consequence of this framework is the pfd sector. Unlike standard
extra dark radiation, which is usually assumed to free-stream and
hence to develop non-zero anisotropic
stress~\cite{Bashinsky2004,Baumann2016,Follin2015}, the PTUM
proxy-fluid behaves as a relativistic \emph{perfect} fluid with
\begin{equation}
w = \tfrac{1}{3},\qquad
c_s^{2} = \tfrac{1}{3},\qquad
\sigma = 0\,.
\label{eq:pfd_signature}
\end{equation}
Free-streaming dark radiation generates a characteristic phase shift
of the CMB acoustic peaks through its anisotropic
stress~\cite{Bashinsky2004,Baumann2016,Follin2015}; the pfd
component, having $\sigma=0$ in the effective linear fluid description,
produces \emph{no such shift}. This absence would constitute a
distinctive discriminator: if additional relativistic degrees of freedom
are detected in CMB data, the presence or absence of the associated
phase shift separates free-streaming from fluid dark radiation. As we
show in Section~\ref{sec:forecast}, CMB-S4 may make this distinction
at high significance for $\Delta N_{\rm fluid}\gtrsim 0.20$,
depending on the realized pfd abundance, making PTUM a concretely
testable cosmological framework.

The aim of the present paper is therefore broader than a
phenomenological dark-radiation study. We pursue three linked goals:

\begin{enumerate}[label=(\roman*)]
\item to formulate PTUM as a cosmological creation scenario based on
      a discharge event, establishing the theoretical framework and
      its effective field theory;
\item to show that the resulting effective dark-sector structure is
      internally consistent within the tested parameter regime, with
      perturbation stability, EFT validity, and compatibility with
      BBN and isocurvature constraints;
\item to test the observable consequences of the model against
      current cosmological data through a Bayesian benchmark analysis,
      profile likelihood, and dataset-split diagnostics.
\end{enumerate}

\paragraph{Scope and claim boundaries.}
Throughout this work, PTUM is treated as an effective cosmological framework,
not as a completed first-principles theory of the origin of the Universe. The
model does not derive the metastable pre-discharge background, does not solve
the cosmological-constant hierarchy by itself, and does not specify a UV
completion. Its present claims are narrower: it defines a threshold-triggered
phase-transition scenario, implements the associated proxy-fluid perturbations,
and identifies an observational discriminator between fluid and free-streaming
dark radiation.


To this end, we proceed in two steps. First, we construct the
effective field theory of PTUM, derive its equations of motion, study
the primordial creation sector---including exact-background
inflationary dynamics---and identify the parameter region in which the
late-time dark-energy sector remains observationally viable. Second,
we implement the proxy-fluid perturbations in the \texttt{CLASS}
Boltzmann code~\cite{Blas2011,Lesgourgues2011}, compute the
cosmological signatures, and constrain the model using a Bayesian benchmark
analysis with \texttt{Cobaya}~\cite{Torrado2021}. The resulting
constraints are cross-checked through profile likelihood extraction
and dataset-split analysis, and confronted with BBN, isocurvature, and
linear stability requirements.

We find that PTUM is observationally viable: the pfd abundance is
bounded to $\Delta N_{\rm fluid}<0.31$ at 95\% credibility from the
full data combination, model comparison gives
$\Delta\mathrm{AIC}=2.26$ relative to $\Lambda$CDM, and all
consistency checks are satisfied. The Hubble constant shifts upward by
$\sim\!0.8\;\mathrm{km\,s^{-1}\,Mpc^{-1}}$ relative to $\Lambda$CDM,
partially alleviating the Hubble tension without being designed to do
so. The discharge energy scale is fixed to
$V_0^{1/4}\sim 10^{15}$--$10^{16}\;\mathrm{GeV}$ by the observed
scalar amplitude, and the creation sector is compatible with current
bounds on the scalar spectral index and the tensor-to-scalar ratio.

The paper is organized as follows.
Section~\ref{sec:conceptual} summarizes the physical picture of PTUM.
Section~\ref{sec:eft} presents the effective field theory,
equations of motion, and symmetry structure.
Section~\ref{sec:primordial} studies the primordial creation sector.
Section~\ref{sec:late} analyses the late-time dark-energy and
dark-sector structure.
Section~\ref{sec:pert} describes the perturbation framework and
\texttt{CLASS} implementation.
Section~\ref{sec:sigs} presents the cosmological signatures.
Section~\ref{sec:mcmc} gives the Bayesian constraints, profile
likelihood, dataset-split analysis, and model comparison.
Section~\ref{sec:stability} collects the stability and consistency
tests.
We discuss the interpretation, limitations, and future prospects in
Section~\ref{sec:discussion} and conclude in
Section~\ref{sec:conclusion}.

%============================================================
\section{Conceptual overview of PTUM}
\label{sec:conceptual}
%============================================================

PTUM begins with a multiverse-like background containing a scalar
quantum fuel field \(\Phi\). The central idea is that local
accumulation of this field drives a threshold-triggered discharge
process. Once a critical configuration is reached, the effective
mass of a proxy activation mode becomes negative, generating a
tachyonic transition that activates string-like degrees of freedom.

The basic sequence is:

\begin{description}
\item[(a) Fuel accumulation.] The scalar fuel field \(\Phi\) accumulates
in a metastable environment. The effective potential is chosen such
that it is approximately quadratic at small field values but saturates
to a plateau at large field values.

\item[(b) Threshold discharge.] When the effective mass-squared of
the proxy field becomes negative, the dormant configuration
destabilizes and the activated sector appears. This transition plays
the role of the creation trigger.

\item[(c) Residual splitting.] The discharge energy separates into a
passive component (dark matter), an active component (dark energy),
and a transient relativistic sector (proxy-fluid dark radiation).
\end{description}

\begin{figure*}[t]
\centering
\includegraphics[width=0.88\textwidth]{fig_model_flowchart.pdf}
\caption{Conceptual flow of the PTUM scenario: accumulation of the
quantum fuel field, threshold discharge, activation of proxy
string-like modes, and residual splitting into dark matter, dark
energy, and proxy-fluid dark radiation.}
\label{fig:flowchart}
\end{figure*}

This picture gives PTUM a structural unification property: the dark
sector is not postulated as a collection of disconnected components,
but emerges from one creation event.

%============================================================
\section{Effective field theory and equations of motion}
\label{sec:eft}
%============================================================

\subsection{Field content}

The minimal PTUM field content consists of four sectors:

\begin{itemize}
\item \(\Phi\): quantum fuel field,
\item \(\psi\): activated string-like proxy mode,
\item \(\chi\): dark-matter field,
\item \(\Phi_{\rm res}\): residual energy field (dark energy).
\end{itemize}

\begin{table}[t]
\centering
\small
\caption{Minimal PTUM field and sector registry.}
\label{tab:ptum_field_registry}
\begin{tabular}{@{}>{\raggedright\arraybackslash}p{0.20\linewidth}>{\raggedright\arraybackslash}p{0.26\linewidth}>{\raggedright\arraybackslash}p{0.40\linewidth}@{}}
\toprule
Symbol & Sector & Effective role \\
\midrule
$\Phi$ & fuel field & accumulates and controls threshold crossing \\
$\psi$ & proxy activation mode & tachyonic activation; proxy-fluid radiation \\
$\chi$ & dark-matter candidate & passive post-discharge matter-like sector \\
$\Phi_{\rm res}$ & residual field & near-$\Lambda$ dark-energy-like sector \\
$\Delta N_{\rm fluid}$ & pfd abundance & maps proxy radiation to CMB/BBN observables \\
\bottomrule
\end{tabular}
\end{table}

\subsection{PTUM Lagrangian}

The total effective Lagrangian is
\begin{equation}
\mathcal{L}_{\rm PTUM}
=\mathcal{L}_{\rm fuel}
+\mathcal{L}_{\rm activation}
+\mathcal{L}_{\rm DM}
+\mathcal{L}_{\rm DE}.
\end{equation}

The fuel sector is
\begin{equation}
\mathcal{L}_{\rm fuel}
=\tfrac{1}{2}(\partial_\mu\Phi)^2
-V_{\rm eff}(\Phi),
\end{equation}
with
\begin{equation}
\begin{split}
V_{\rm eff}(\Phi)
&=V_0\!\left(1-e^{-\lambda\Phi^2}\right) \\
&\quad +\frac{1}{64\pi^2}
\sum_i n_i\,m_{{\rm eff},i}^4(\Phi)
\left[\ln\frac{m_{{\rm eff},i}^2(\Phi)}{\mu_R^2}
-\tfrac{3}{2}\right].
\end{split}
\label{eq:veff}
\end{equation}
The leading tree-level part is plateau-like; the second term
represents Coleman--Weinberg corrections~\cite{ColemanWeinberg}.

The activation sector is
\begin{equation}
\mathcal{L}_{\rm activation}
=\tfrac{1}{2}(\partial_\mu\psi)^2
-\frac{m^2}{2}\psi^2\!\left(1-\frac{\alpha\Phi^2}{E_{\rm th}^2}\right)
-\frac{\lambda_\psi}{4}\psi^4,
\label{eq:lact}
\end{equation}
with effective mass
\begin{equation}
m_{\rm eff}^2(\Phi)
=m^2\!\left(1-\frac{\alpha\Phi^2}{E_{\rm th}^2}\right).
\end{equation}
Equivalently, the effective discharge trigger can be written as
\begin{equation}
\boxed{
 m_{\rm eff}^2(\Phi)<0
 \quad\Longleftrightarrow\quad
 \Phi>\Phi_c\equiv \frac{E_{\rm th}}{\sqrt{\alpha}} .
}
\label{eq:ptum_threshold_box}
\end{equation}
This condition defines the transition from the dormant proxy phase to the
activated proxy phase. It should be interpreted as an effective threshold
criterion, not as a microscopic derivation of the pre-discharge background.
Here, $E_{\rm th}$ is the dimensionful discharge threshold scale (naturally related to the creation energy $V_0^{1/4}$), and $\alpha$ is a dimensionless coupling constant. The tachyonic instability is triggered when the fuel field crosses the critical threshold $\Phi = E_{\rm th}/\sqrt{\alpha}$. The transition timescale is assumed to be much shorter than the local Hubble time ($H^{-1}$), allowing the discharge to be treated as a nearly instantaneous phase transition at the background level that partitions the available energy into the residual dark sector and standard radiation. For \(m_{\rm eff}^2>0\) the proxy mode is dormant; for \(m_{\rm eff}^2<0\) it is activated by the tachyonic instability.

The dark-matter sector is
\begin{equation}
\mathcal{L}_{\rm DM}
=\tfrac{1}{2}(\partial_\mu\chi)^2
-\tfrac{1}{2}m_\chi^2\chi^2
-g_{\rm DM}\chi^2\Phi_{\rm res}^2,
\end{equation}
and the dark-energy sector is
\begin{equation}
\mathcal{L}_{\rm DE}
=\tfrac{1}{2}(\partial_\mu\Phi_{\rm res})^2
-W(\Phi_{\rm res}),
\end{equation}
with
\begin{equation}
\begin{split}
W(\Phi_{\rm res})
&=W_0\,e^{-\kappa\Phi_{\rm res}^2}
\!\left(1-e^{-\lambda_{\rm res}\Phi_{\rm res}^2}\right) \\
&\quad +\Lambda_{\rm NP}^4\,e^{-S_E}
\cos\!\left(\frac{\Phi_{\rm res}}{f_{\rm res}}\right).
\end{split}
\label{eq:wpot}
\end{equation}

\subsection{Equations of motion}

The field equations in a flat FRW universe,
\(ds^2=-dt^2+a^2(t)d\mathbf{x}^2\), read
\begin{align}
\ddot\Phi+3H\dot\Phi
+\frac{\partial V_{\rm eff}}{\partial\Phi}
+\frac{\alpha m^2}{E_{\rm th}^2}\Phi\psi^2
&=0,
\label{eq:frw1}
\\
\ddot\psi+3H\dot\psi
+m^2\!\left(1-\frac{\alpha\Phi^2}{E_{\rm th}^2}\right)\psi
+\lambda_\psi\psi^3
&=0,
\label{eq:frw2}
\\
\ddot\chi+3H\dot\chi+m_\chi^2\chi
+2g_{\rm DM}\chi\Phi_{\rm res}^2
&=0,
\label{eq:frw3}
\\
\ddot\Phi_{\rm res}+3H\dot\Phi_{\rm res}
+\frac{\partial W}{\partial\Phi_{\rm res}}
+2g_{\rm DM}\chi^2\Phi_{\rm res}
&=0.
\label{eq:frw4}
\end{align}

\subsection{Quartic attractor and the proxy-fluid sector}

A key result is that coherent oscillations of \(\psi\) in a quartic
potential behave effectively as radiation~\cite{Turner1983}:
\begin{equation}
w_\psi=\tfrac{1}{3},\qquad \rho_\psi\propto a^{-4}.
\end{equation}
However, unlike free-streaming radiation, the coherent scalar
condensate does not develop a full Boltzmann hierarchy and therefore
contributes no anisotropic stress:
\begin{equation}
\sigma_\psi=0.
\end{equation}
In the effective linear fluid closure used below, this sector is therefore
specified by
\begin{equation}
\boxed{
 w_\psi=\frac{1}{3},\qquad
 c_{s,\psi}^2=\frac{1}{3},\qquad
 \sigma_\psi=0 .
}
\label{eq:pfd_closure_box}
\end{equation}
The third condition is the key discriminator: the proxy-fluid component does not
carry a free-streaming hierarchy and therefore does not generate the
neutrino-like CMB acoustic phase shift associated with anisotropic stress.
This defines the proxy-fluid dark-radiation sector. The absence of
anisotropic stress is the defining observational signature of the
pfd component: free-streaming dark radiation with the same energy
density would produce a characteristic CMB acoustic-peak phase
shift~\cite{Bashinsky2004,Follin2015}, whereas pfd does not.

The difference between the PTUM proxy-fluid component and a conventional
free-streaming relativistic species is summarized schematically in
Fig.~\ref{fig:ptum_proxy_fluid_vs_free_streaming}. The two cases may contribute
a similar amount of extra radiation density, but they differ at the perturbative
level: the proxy-fluid sector has no effective anisotropic stress in the linear
fluid description, while a free-streaming component sources a neutrino-like
acoustic phase shift.

\begin{figure*}[p]
\centering
\includegraphics[
width=0.98\textwidth,
height=0.62\textheight,
keepaspectratio
]{figures/PTUM_F3_proxy_fluid_vs_free_streaming_dark_radiation.pdf}
\caption{
Proxy-fluid versus free-streaming dark radiation. In the effective linear
fluid description, the PTUM proxy component has
$w_\psi=1/3$, $c_{s,\psi}^2=1/3$, and $\sigma_\psi=0$, and therefore does
not induce the neutrino-like CMB acoustic phase shift characteristic of
free-streaming relativistic species. The key discriminator is whether an
extra relativistic component produces a measurable acoustic phase shift, not
only whether it contributes to the radiation density.
}
\label{fig:ptum_proxy_fluid_vs_free_streaming}
\end{figure*}


\subsection{Operational stability and observable closure of the discharge sector}
\label{subsec:ptum_operational_stability}

The preceding subsections define PTUM as a threshold-triggered discharge
mechanism: the fuel field $\Phi$ accumulates, the proxy field $\psi$ becomes
tachyonically active when its effective mass squared changes sign, and the
post-discharge energy is partitioned into dark matter, residual dark energy,
and proxy-fluid dark radiation.  The present subsection records the minimal
internal stability conditions required for this effective description to be
well defined.  No additional field, sector, or cross-model structure is
introduced here.

The discharge threshold is determined by the proxy-field effective mass,
\begin{equation}
m_{\rm eff}^2(\Phi)
=
m^2
\left(
1-\frac{\alpha\Phi^2}{E_{\rm th}^2}
\right),
\label{eq:PTUM_meff_operational}
\end{equation}
so that the critical field value is
\begin{equation}
\Phi_c
=
\frac{E_{\rm th}}{\sqrt{\alpha}} .
\label{eq:PTUM_critical_phi_operational}
\end{equation}
For $|\Phi|<\Phi_c$ the proxy sector is dormant, while for
$|\Phi|>\Phi_c$ the tachyonic sign change activates the discharge channel.
The activation potential may be written as
\begin{equation}
V_{\rm act}(\psi;\Phi)
=
\frac{1}{2}m_{\rm eff}^2(\Phi)\psi^2
+
\frac{\lambda_\psi}{4}\psi^4 .
\label{eq:PTUM_activation_potential_operational}
\end{equation}
When $m_{\rm eff}^2<0$, the post-activation minimum is
\begin{equation}
\psi_0^2
=
-\frac{m_{\rm eff}^2(\Phi)}{\lambda_\psi},
\label{eq:PTUM_proxy_minimum_operational}
\end{equation}
and local stability of the activated state requires
\begin{equation}
\left.
\frac{\partial^2 V_{\rm act}}{\partial \psi^2}
\right|_{\psi=\psi_0}
=
-2m_{\rm eff}^2(\Phi)
>
0 .
\label{eq:PTUM_potential_stability_operational}
\end{equation}
Thus the tachyonic sign change is not treated as an endpoint instability;
it is a trigger for a transition to a locally stable post-discharge state
inside the effective theory.

Averaged over coherent oscillations in the quartic potential, the proxy sector
behaves as radiation,
\begin{equation}
w_\psi=\frac{1}{3},
\qquad
\rho_\psi\propto a^{-4}.
\label{eq:PTUM_quartic_fluid_operational}
\end{equation}
In the effective linear closure used in the Boltzmann implementation, the
proxy component is represented as a perfect relativistic fluid,
\begin{equation}
c_{s,\psi}^2=\frac{1}{3},
\qquad
\sigma_\psi=0 .
\label{eq:PTUM_proxy_signature_operational}
\end{equation}
This last condition is the key observable distinction between PTUM proxy-fluid
dark radiation and free-streaming dark radiation.

Combining the internal and observable requirements, the operational PTUM
closure used in the remainder of this paper is
\begin{equation}
\boxed{
\begin{gathered}
\text{activation: } m_{\rm eff}^2(\Phi)<0,\\[2pt]
\text{stability: } -2m_{\rm eff}^2(\Phi)>0,\\[2pt]
\text{fluid closure: } w_\psi=1/3,\; c_{s,\psi}^2=1/3,\; \sigma_\psi=0 .
\end{gathered}}
\label{eq:ptum_operational_closure}
\end{equation}
The first two lines are internal EFT consistency requirements.  The final line
is the perturbative input tested through the CMB acoustic-phase discriminator
and the late-time expansion analysis below.

\subsection{Symmetry structure}

The symmetry structure of the effective theory includes:

\begin{itemize}
\item an exact \(\mathbb{Z}_2\) symmetry \(\Phi\to-\Phi\),
\item a spontaneously broken \(\mathbb{Z}_2\) symmetry
      \(\psi\to-\psi\),
\item an exact \(\mathbb{Z}_2\) symmetry \(\chi\to-\chi\) ensuring
      dark-matter stability,
\item an approximate shift symmetry
      \(\Phi_{\rm res}\to\Phi_{\rm res}+c\) protecting the
      dark-energy sector.
\end{itemize}

\subsection{UV perspective and effective field theory scope}
\label{sec:uv_perspective}

The PTUM Lagrangian is constructed as a low-energy effective field
theory valid at energy scales $E \ll M_{\rm Pl}$. We do not claim
that PTUM constitutes a UV-complete theory; rather, it parametrises
the phenomenological consequences that a broad class of UV
completions would produce at cosmological scales. Several structural
features of PTUM find natural counterparts in well-studied
high-energy frameworks, which we briefly discuss below. The UV comments in
this subsection are therefore interpretive guideposts rather than required
assumptions of the model. PTUM should remain testable at the EFT and
Boltzmann-implementation level even if no particular string-theoretic,
swampland, or higher-dimensional embedding is adopted.

\noindent\textbf{Tachyonic transition and string theory.}
The activation mechanism through which the fuel field $\Phi$ triggers
the proxy mode $\psi$ via a tachyonic mass term bears a close
structural resemblance to the open-string tachyon condensation that
arises in D-brane--anti-D-brane systems
\cite{Sen2002,Dvali1999,KKLMMT2003}. In the Dvali--Tye scenario, the
inter-brane separation plays the role of an inflaton, and brane
annihilation ends inflation through tachyon condensation. In PTUM,
$\Phi$ plays an analogous role to the inter-brane modulus, while
$\psi$ corresponds to the open-string tachyon that becomes active
once $\Phi$ crosses a critical threshold. We emphasise that this
correspondence is at the level of structural analogy rather than
rigorous derivation; a full embedding into a specific string
compactification is left for future work.

\noindent\textbf{Swampland considerations.}
The fuel potential $V(\Phi) = V_0(1 - e^{-\lambda\Phi^2})$ satisfies
the de Sitter Swampland Conjecture \cite{Obied2018,Ooguri2019} in
the large-field regime where $|V'|/V \gtrsim \mathcal{O}(1)$, while
the inflationary slow-roll regime ($\epsilon, |\eta| \ll 1$) is
maintained for a finite number of $e$-folds. The refined de Sitter
conjecture, which allows $\min(\nabla_i \nabla_j V) \leq
-c'\,V/M_{\rm Pl}^2$, is satisfied throughout the slow-roll phase
for $c' \sim \mathcal{O}(\eta)$. A systematic analysis of the full
PTUM potential landscape against swampland criteria is deferred to
future work.

\noindent\textbf{Field excursions and shift symmetry.}
For small values of the coupling $\lambda \lesssim 0.05$, the fuel
field undergoes super-Planckian excursions $\Delta\Phi > M_{\rm Pl}$
during the inflationary phase. This is a generic feature of
large-field inflation models and is not specific to PTUM. Such
excursions can be protected against dangerous higher-dimensional
operators by an approximate shift symmetry $\Phi \to \Phi + c$,
which is broken only by the exponential potential. This is the same
protective mechanism invoked in natural inflation
\cite{Freese1990} and axion monodromy \cite{Silverstein2008} models.

\noindent\textbf{EFT validity summary.}
As demonstrated in Section~\ref{sec:stability}, all PTUM energy
scales satisfy $H_* \ll M_{\rm Pl}$, all couplings remain
perturbative, and the theory is free from ghost, gradient, and
tachyonic instabilities. The EFT is therefore self-consistent within
its domain of applicability. The question of which specific UV
completion gives rise to the PTUM effective Lagrangian is an
important but separate question that does not affect the
observational predictions derived in this work.

%============================================================
\section{Primordial consistency of the creation sector}
\label{sec:primordial}
%============================================================

The primordial consistency of PTUM is tested by treating the fuel
potential as the effective generator of an early accelerated phase.
We ask whether the creation sector can be normalized consistently
to the observed primordial scalar amplitude and whether its
predictions for \(n_s\) and \(r\) fall within the
Planck/BICEP-allowed region.

\subsection{Slow-roll analysis}

At tree level, the fuel potential is
\begin{equation}
V(\Phi)=V_0\!\left(1-e^{-\lambda\Phi^2}\right).
\label{eq:fuelpot}
\end{equation}
This is a plateau-type potential: approximately quadratic at small
\(\Phi\) and saturating to a constant at large \(\Phi\). Such
plateau structures are strongly favored by CMB data~\cite{BICEPKeck2021}.

We define the potential slow-roll parameters
\begin{equation}
\epsilon_V=\tfrac{1}{2}\!\left(\frac{V_{,\Phi}}{V}\right)^2,\qquad
\eta_V=\frac{V_{,\Phi\Phi}}{V},
\end{equation}
and the standard approximate observables
\begin{equation}
n_s\simeq1-6\epsilon_V+2\eta_V,\qquad r\simeq16\epsilon_V.
\end{equation}
A scan over \(\lambda\) shows that for \(\lambda\gtrsim0.02\) the
PTUM creation sector naturally predicts \(n_s\sim0.957\)--\(0.961\)
and \(r\ll0.036\), compatible with current
constraints~\cite{BICEPKeck2021,Planck2018params}.

\subsection{Non-Gaussianity}

Because the creation sector is a canonical single-field slow-roll
model, the non-Gaussianity parameters satisfy the Maldacena
consistency relation~\cite{Maldacena2003}:
\begin{equation}
f_{\rm NL}^{\rm local}
=\tfrac{5}{12}(1-n_s)\sim0.017\text{--}0.018,\qquad
f_{\rm NL}^{\rm equil}\sim10^{-3}.
\end{equation}
These values are well within current Planck bounds
(\(|f_{\rm NL}^{\rm local}|<5.1\),
\(|f_{\rm NL}^{\rm equil}|<47\) at 95\%~CL~\cite{PlanckNG})
by factors of order \(10^2\)--\(10^3\), and lie far below the
sensitivity of any planned experiment. This constitutes a clean
null prediction of PTUM.

\subsection{Exact-background consistency}

To verify the slow-roll treatment we performed an exact-background
integration using the bounded variable
\begin{equation}
X\equiv\frac{\dot\Phi}{\sqrt{6}\,H},\qquad \epsilon_H=3X^2.
\end{equation}
The exact background equations in terms of \(N=\ln a\) are
\begin{align}
\frac{d\Phi}{dN}&=\sqrt{6}\,X,\\[4pt]
\frac{dX}{dN}&=(1-X^2)\!\left(\sqrt{\tfrac{3}{2}}\,s(\Phi)-3X\right),
\end{align}
where \(s(\Phi)\equiv-V_{,\Phi}/V\). For each \(\lambda\) we searched
for an initial field value \(\Phi_i\) yielding \(N_{\rm end}\sim60\)
and evaluated the observables 55 e-folds before the end. Results are
given in Table~\ref{tab:exactscan}.

\begin{table*}
\caption{Exact-background primordial scan of the PTUM creation
sector. The exact and slow-roll predictions agree to within
\(\Delta n_s\simeq5\times10^{-4}\) and
\(|\Delta r|/r\simeq1.5\%\).}
\label{tab:exactscan}
\centering
\begin{tabular*}{\textwidth}{@{\extracolsep{\fill}}ccccccccc}
\hline\noalign{\smallskip}
\(\lambda\) & \(\Phi_i\) & \(N_{\rm end}\) & \(\Phi_*\) &
\(\epsilon_*\) & \(n_s^{\rm ex}\) & \(n_s^{\rm SR}\) &
\(r^{\rm ex}\) & Status \\
\noalign{\smallskip}\hline\noalign{\smallskip}
0.020 & 10.722 & 60.14 & 10.465 & \(1.37\times10^{-3}\) &
0.9581 & 0.9576 & 0.02195 & Compat. \\
0.050 &  8.325 & 59.96 &  8.185 & \(4.37\times10^{-4}\) &
0.9565 & 0.9559 & 0.00699 & Boundary \\
0.100 &  6.612 & 59.84 &  6.527 & \(1.72\times10^{-4}\) &
0.9564 & 0.9559 & 0.00276 & Boundary \\
0.200 &  5.137 & 60.38 &  5.078 & \(6.82\times10^{-5}\) &
0.9570 & 0.9564 & 0.00109 & Compat. \\
0.500 &  3.579 & 59.52 &  3.551 & \(2.08\times10^{-5}\) &
0.9579 & 0.9574 & 0.000333 & Compat. \\
1.000 &  2.691 & 59.32 &  2.673 & \(8.76\times10^{-6}\) &
0.9586 & 0.9580 & 0.000140 & Compat. \\
2.000 &  2.012 & 60.99 &  1.996 & \(3.77\times10^{-6}\) &
0.9591 & 0.9586 & 0.000060 & Compat. \\
5.000 &  1.350 & 58.88 &  1.344 & \(1.28\times10^{-6}\) &
0.9597 & 0.9592 & 0.000020 & Compat. \\
\noalign{\smallskip}\hline
\end{tabular*}
\end{table*}

\begin{figure}[t]
\centering
\includegraphics[width=0.95\columnwidth]{fig_ns_r_plane_exact.pdf}
\caption{PTUM predictions in the \(n_s\)--\(r\) plane using the
exact-background calculation. Blue bands indicate the
Planck-preferred region; the green shading marks the BICEP/Keck
upper bound on \(r\). For \(\lambda\gtrsim0.02\), PTUM lies in or
very near the observationally allowed region.}
\label{fig:nsr_exact}
\end{figure}

\begin{figure*}[t]
\centering
\includegraphics[width=0.95\textwidth]{fig_exact_vs_SR_all_lambda.pdf}
\caption{Comparison of exact-background and slow-roll primordial
predictions across the full \(\lambda\) scan. The slow-roll
approximation underestimates \(n_s\) by \(\sim5\times10^{-4}\) and
overestimates \(r\) by \(\sim1.4\%\), confirming its quantitative
reliability for the PTUM creation sector.}
\label{fig:exactvssr}
\end{figure*}

\subsection{Primordial normalization and creation scale}

Matching the observed scalar amplitude
\(A_s\simeq2.1\times10^{-9}\) via
\begin{equation}
A_s\simeq\frac{V_*}{24\pi^2 M_{\rm Pl}^4\,\epsilon_*},
\end{equation}
gives
\begin{equation}
V_0^{1/4}\sim2\times10^{15}\text{--}1.3\times10^{16}\ {\rm GeV},
\qquad
H_*\sim10^{12}\text{--}10^{13}\ {\rm GeV}.
\end{equation}

\begin{figure}[t]
\centering
\includegraphics[width=0.95\columnwidth]{fig_PTUM_creation_scale.pdf}
\caption{Characteristic PTUM creation scales inferred from the
observed scalar amplitude. The PTUM creation event lies in a
GUT-adjacent energy band, well below the Planck scale.}
\label{fig:creation}
\end{figure}

This places the PTUM discharge event in a physically meaningful
GUT-adjacent regime. The EFT hierarchy is satisfied throughout the benchmark regime:
\(H_*/M_{\rm Pl}\sim4\times10^{-6}\ll1\) and
\(V_0/M_{\rm Pl}^4\sim10^{-10}\ll1\).

%============================================================
\section{Late-time dark sector viability}
\label{sec:late}
%============================================================

\subsection{Dark-energy sector}

The PTUM dark-energy potential~\eqref{eq:wpot} was scanned over the
\((\kappa,\lambda_{\rm res})\) plane. In 89.6\% of the scanned
parameter space there exists a field value giving
\(1+w_{\rm DE}<0.01\), and in 95.1\% there exists one giving
\(1+w_{\rm DE}<0.05\).

\begin{figure}[t]
\centering
\includegraphics[width=0.95\columnwidth]{fig_DE_viable_region.pdf}
\caption{Viable-region scan of the PTUM dark-energy sector.
Left: best achievable \(\log_{10}(1+w_{\rm DE})\) in the
\((\kappa,\lambda_{\rm res})\) plane. Right: corresponding optimal
field value \(\Phi_{\rm res}\). The wide viable region shows that
near-\(\Lambda\) behavior is not confined to a fine-tuned point.}
\label{fig:DEviable}
\end{figure}

\subsection{Dark-sector evolution function}

Using the CLASS best-fit background solution we define the
dark-energy fraction
\begin{equation}
\gamma(z)
\equiv
\frac{\rho_{\rm DE}(z)}{\rho_{\rm DM}(z)+\rho_{\rm DE}(z)}.
\end{equation}
The numerical evolution is accurately described by
\begin{equation}
\gamma(z)
=\frac{1}{1+c\,(1+z)^3},
\qquad
c=\frac{\rho_{{\rm DM},0}}{\rho_{{\rm DE},0}},
\label{eq:gamma}
\end{equation}
with best-fit value \(c=0.3854\), yielding \(\gamma(0)=0.722\).
The analytic form is verified against the full CLASS output at the
level of numerical precision (residuals
\(|\Delta\gamma/\gamma|<10^{-14}\) across \(z=0\)--\(5\)).

\begin{figure}[t]
\centering
\includegraphics[width=0.85\columnwidth]{fig_gamma_z_numerical_vs_analytic.pdf}
\caption{The PTUM dark-sector evolution function \(\gamma(z)\)
compared with the CLASS numerical solution. The analytic
form~\eqref{eq:gamma} matches the Boltzmann output within machine
precision across the full redshift range shown.}
\label{fig:gamma}
\end{figure}

\subsection{Post-creation energy partition}

The discharge energy distributes into four sectors: SM radiation
(dominant), dark matter \(\chi\), pfd \(\psi\), and residual dark
energy \(\Phi_{\rm res}\). The pfd fraction relative to radiation is
\begin{equation}
\frac{\rho_\psi}{\rho_{\rm rad}}
=\frac{\Delta N_{\rm fluid}}{N_{\rm eff}^{\rm std}}
\sim3\text{--}11\%
\quad(0.10\le\Delta N_{\rm fluid}\le0.31),
\end{equation}
requiring an effective coupling \(\mu/\Lambda\sim0.18\text{--}0.33\),
which is natural (\(\mathcal{O}(0.1)\), no fine-tuning). The
residual dark-energy fraction is
\(V_{\rm res}/V_0\sim10^{-108}\), which re-parameterizes rather
than solves the cosmological constant problem.

\subsection{Proxy-fluid dark radiation}

The pfd sector is parameterized by
\begin{equation}
\Delta N_{\rm fluid}
=\frac{\rho_{\rm pfd}}
{\tfrac{7}{8}\!\left(\tfrac{4}{11}\right)^{4/3}\!\rho_\gamma}.
\end{equation}
Its defining properties are \(w=1/3\), \(c_s^2=1/3\),
and \(\sigma=0\).

%============================================================
\section{Perturbations and numerical implementation}
\label{sec:pert}
%============================================================

We implemented the PTUM proxy-fluid dark-radiation sector in
CLASS~v3.2.0~\cite{Blas2011,Lesgourgues2011}. In synchronous
gauge~\cite{MaBertschinger}, the scalar perturbations obey
\begin{align}
\dot\delta_{\rm pfd}
&=-\tfrac{4}{3}\!\left(\theta_{\rm pfd}+\tfrac{\dot h}{2}\right),
\\[4pt]
\dot\theta_{\rm pfd}
&=\tfrac{k^2}{4}\delta_{\rm pfd}.
\end{align}

\subsection{Proxy-fluid dark radiation: origin of $\sigma_\psi=0$}
\label{sec:pfd_derivation}

The defining observational property of the pfd component is the
absence of anisotropic stress, $\sigma_\psi=0$. Within the effective coherent-field description, this is not introduced
as an independent assumption; it follows from the nature of the source
field and the regime in which it evolves. Here we derive this result from the
equations of motion.

\subsubsection{Effective energy-momentum tensor}

Consider the proxy field $\psi$ oscillating coherently in its quartic
potential $V(\psi)=\frac{1}{4}\lambda_\psi\psi^4$. The canonical
energy-momentum tensor is
\begin{equation}
T_{\mu\nu}^{(\psi)}
=\partial_\mu\psi\,\partial_\nu\psi
-g_{\mu\nu}\Bigl[\tfrac12(\partial_\rho\psi)^{2}
-V(\psi)\Bigr].
\end{equation}
For a spatially homogeneous field $\psi=\psi(t)$ oscillating in the
quartic potential, the spatial average over one oscillation period
gives the effective fluid description. Using the fact that
$\langle\dot\psi^2\rangle=\langle V(\psi)\rangle$ for $w=1/3$
oscillations, the pressure and energy density satisfy
$p_\psi=\rho_\psi/3$ and the anisotropic stress vanishes in the
effective fluid description:
\begin{equation}
\sigma_\psi \equiv 
\frac{\rho_\psi-p_\psi-2p_\perp}{\rho_\psi+p_\psi}=0\,,
\end{equation}
because the energy-momentum tensor of a spatially homogeneous
coherent field is isotropic at the background level.

This is the key physical point: unlike a collection of free-streaming
particles, which carry momentum correlations through their individual
trajectories and thereby generate non-zero $\sigma$, a coherent
condensate does not have independent degrees of freedom below the
coherence scale. The pressure is everywhere $p=\rho/3$ and there are
no viscous or anisotropic contributions from particle collisions.

\subsubsection{Perturbation hierarchy}

At the linear perturbation level we write $\psi=\bar\psi(t)+\delta\psi$.
The perturbed energy-momentum tensor contains
$\delta\rho_\psi$, $\delta p_\psi$, velocity $\theta_\psi$, and shear
$\sigma_\psi$. The conservation equation
$\nabla_\mu T^{\mu\nu(\psi)}=0$ gives two first-order equations
for $\delta_\psi\equiv\delta\rho_\psi/\bar\rho_\psi$ and
$\theta_\psi\equiv ik\delta v_\psi$:
\begin{align}
\dot\delta_\psi
&=-\frac{4}{3}\Bigl(\theta_\psi+\frac{\dot h}{2}\Bigr)\,,
\label{eq:delta_pfd_main}\\[4pt]
\dot\theta_\psi
&=\frac{k^2}{4}\,\delta_\psi\,.
\label{eq:theta_pfd_main}
\end{align}
Higher multipoles ($F_\ell$ with $\ell\ge2$) satisfy
$\dot F_\ell\propto k F_{\ell-1}-(\ell+1)H F_\ell$ and do not source
the Einstein equations for scalar modes. These higher multipoles are
present but they do not contribute to $\sigma_\psi$ in the coherent
oscillation limit because the condensate has no particle-streaming
velocity dispersion. They are therefore omitted from the evolution
equations used in the \texttt{CLASS} implementation.

The conditions under which this approximation holds are:
\begin{itemize}
\item \textbf{Coherence:} The coherence length of the $\psi$
condensate exceeds the Hubble radius, so that the field oscillates as
a single macroscopic mode. This is naturally satisfied when the
field is homogeneous at the scale of the observable patch at the time
of oscillatory onset.
\item \textbf{Non-interacting:} The $\psi$ field has no significant
self-interaction beyond the $\lambda_\psi\psi^4$ potential term.
If $\lambda_\psi\psi^4$ is the only relevant interaction, the
condensate remains collisionless and coherent.
\item \textbf{Subluminal sound horizon:} The physical sound horizon
$c_s t$ of the condensate is sub-Hubble, ensuring that the perfect
fluid approximation is valid.
\end{itemize}

The sound speed follows from the perturbation equations:
$c_s^2=\dot p_\psi/\dot\rho_\psi=1/3$, as expected for a
relativistic coherent oscillator. This value is positive and
constant, confirming gradient stability ($c_s^2>0$) throughout the
evolution.

We emphasise that this is the same physics underlying the standard
treatment of scalar field dark energy (quintessence) and inflation
ensuring $c_s^2=1$ for the inflaton. The pfd case is distinguished
only by the fact that the condensate is in the $w=1/3$ regime rather
than $w\simeq-1$ or $w\simeq0$.

The fluid is therefore a \emph{relativistic perfect fluid} with
$c_s^2=1/3$ and $\sigma=0$, and its perturbation equations reduce
to the two coupled first-order equations (\ref{eq:delta_pfd_main})--(\ref{eq:theta_pfd_main}), with no shear term and no higher-multipole hierarchy. 

The implementation supplies adiabatic initial conditions (consistent
with the single-field creation mechanism, which yields purely
adiabatic perturbations to leading order in slow-roll) and couples
the pfd fluid consistently to the Einstein source terms. Exact
\(\Lambda\)CDM recovery for \(\Delta N_{\rm fluid}=0\) was verified.

A technically important issue was the state-vector layout. An early
implementation inserted pfd variables inside the ultra-relativistic
hierarchy, breaking the contiguous index assumptions used by CLASS
and producing singular-matrix failures. The final implementation
places pfd in an independent fluid block, fully resolving the
instability.

%============================================================
\section{Cosmological signatures}
\label{sec:sigs}
%============================================================

\subsection{CMB residuals}

Raw CMB TT and matter-power residuals were computed for
\(\Delta N_{\rm fluid}=0.01\), \(0.02\), \(0.05\), \(0.10\),
\(0.20\), keeping the standard cosmological parameters fixed.

\begin{figure*}[t]
\centering
\includegraphics[width=0.95\textwidth]{fig_multi_dN_residuals_raw.pdf}
\caption{Raw CMB TT residuals (left) and matter-power residuals
(right) for PTUM proxy-fluid dark radiation over
\(\Delta N_{\rm fluid}=0.01\)--\(0.20\). The response scales
approximately linearly with \(\Delta N_{\rm fluid}\).}
\label{fig:multiDN}
\end{figure*}

\subsection{Lensing potential}

The CMB lensing potential power spectrum \(C_L^{\phi\phi}\) is
suppressed by pfd through its effect on late-time structure growth.
At \(L=500\) the fractional suppression is
\begin{equation}
\begin{aligned}
\frac{\Delta C_L^{\phi\phi}}{C_L^{\phi\phi}}
&\simeq -1.46\%\quad(\Delta N=0.10),\\
&\simeq -4.67\%\quad(\Delta N=0.31).
\end{aligned}
\end{equation}
The suppression scales approximately linearly with
\(\Delta N_{\rm fluid}\) and is detectable with CMB-S4 lensing
reconstruction for \(\Delta N_{\rm fluid}\gtrsim0.10\).

\begin{figure}[t]
\centering
\includegraphics[width=0.95\columnwidth]{fig_lensing_potential_PTUM.pdf}
\caption{Fractional residuals of the CMB lensing potential power
spectrum \(C_L^{\phi\phi}\) for PTUM relative to \(\Lambda\)CDM.
The suppression grows with \(\Delta N_{\rm fluid}\) and reaches
\(\sim5\%\) at \(L\sim1000\) for \(\Delta N=0.31\).}
\label{fig:lensing_pp}
\end{figure}

\subsection{Fluid versus free-streaming dark radiation}

The key observable distinction between pfd and standard
neutrino-like dark radiation is the absence of anisotropic stress
in the pfd case. Free-streaming dark radiation with extra energy
fraction \(\Delta N_{\rm eff}\) induces a CMB acoustic-peak phase
shift
\begin{equation}
\begin{split}
\delta\phi_{\rm peak}
&\approx 0.19\pi\,R_{\rm extra}, \\
R_{\rm extra}
&=\frac{\Delta N_{\rm fluid}\,\tfrac{7}{8}
\!\left(\tfrac{4}{11}\right)^{4/3}}
{\sum_i\Omega_{i,{\rm rad}}},
\end{split}
\end{equation}
whereas pfd with \(\sigma=0\) produces no such shift.

The observational distinction between a proxy-fluid component and a free-streaming relativistic species can be expressed directly in terms of the acoustic phase of the CMB temperature spectrum. For a fixed additional radiation abundance, a free-streaming component shifts the acoustic peaks through anisotropic-stress effects, whereas the PTUM proxy-fluid component preserves the standard peak phase within the effective linear fluid description. This diagnostic is illustrated schematically in Fig.~\ref{fig:ptum_cmb_phase_shift_diagnostic}.

\begin{figure*}[t!]
\centering
\includegraphics[
width=0.92\textwidth,
height=0.52\textheight,
keepaspectratio
]{figures/PTUM_F4_CMB_phase_shift_diagnostic.pdf}
\caption{
CMB acoustic phase-shift diagnostic. A free-streaming relativistic component
generates anisotropic stress and produces a neutrino-like displacement of the
acoustic peak structure relative to the baseline spectrum. In contrast, the
PTUM proxy-fluid component has $\sigma_\psi=0$ within the effective linear
fluid approximation and therefore does not induce the same phase shift. The
curves shown here are schematic diagnostics rather than numerical CLASS
outputs.
}
\label{fig:ptum_cmb_phase_shift_diagnostic}
\end{figure*}


The predicted phase shifts are listed in
Table~\ref{tab:phaseshift}.

At fixed cosmological parameters, a systematic comparison of TT
and EE power spectra shows that the absence of the sigma-induced
phase shift in the pfd case leads to smaller spectral distortions
than in the free-streaming case with the same \(\Delta N\). The
difference is dominated by the EE polarization spectrum and grows
rapidly with \(\Delta N_{\rm fluid}\), indicating that current and
future high-resolution CMB data should be able to discriminate
between the two scenarios. This distinction is most directly
accessible through the CMB acoustic-peak phase, which serves as
the primary observational discriminator of pfd.

\begin{table}
\caption{Phase shift between fluid (pfd) and free-streaming dark
radiation for representative values of \(\Delta N_{\rm fluid}\),
together with the detection significance for CMB-S4 and CMB-HD.}
\label{tab:phaseshift}
\centering
\begin{tabular}{cccc}
\hline\noalign{\smallskip}
\(\Delta N_{\rm fluid}\) & Phase shift & CMB-S4 & CMB-HD \\
\noalign{\smallskip}\hline\noalign{\smallskip}
0.05 & \(0.23^\circ\) & no (0.8\(\sigma\)) & marginal (1.5\(\sigma\)) \\
0.10 & \(0.45^\circ\) & no (1.5\(\sigma\)) & detect (3.0\(\sigma\)) \\
0.20 & \(0.89^\circ\) & marginal (3.0\(\sigma\)) & detect (6.0\(\sigma\)) \\
0.31 & \(1.40^\circ\) & detect (4.5\(\sigma\)) & detect (9.2\(\sigma\)) \\
\noalign{\smallskip}\hline

\end{tabular}
\end{table}

\subsection{Decision-level observational tests}
\label{sec:decision_tests}

The proxy-fluid interpretation should be evaluated through a small number of
separable observational outcomes.  This is important because PTUM is not a
single yes-or-no deformation of \(\Lambda\)CDM: its background expansion,
radiation abundance, perturbative closure, and phase-shift signature probe
different parts of the effective construction.  Table~\ref{tab:decision_tests}
collects the corresponding decision tests.

\begin{table*}[t]
\caption{Decision-level observational tests for the PTUM proxy-fluid sector.
The table separates viability tests, discriminators, and direct failure modes.}
\label{tab:decision_tests}
\centering
\small
\begin{tabular}{p{0.20\textwidth}p{0.24\textwidth}p{0.23\textwidth}p{0.19\textwidth}}
\hline\noalign{\smallskip}
Test & PTUM expectation & Distinguishing information & Failure mode \\
\noalign{\smallskip}\hline\noalign{\smallskip}
Radiation abundance & A bounded positive or vanishing
\(\Delta N_{\rm fluid}\) & Determines whether the proxy-fluid channel is
cosmologically relevant & No allowed region after CMB, BBN, BAO, SNe, and lensing \\
Acoustic phase & No neutrino-like phase shift from the proxy component & Separates fluid radiation from free-streaming radiation & Extra radiation is detected with a free-streaming phase shift \\
Perturbative closure & \(w_\psi=1/3\), \(c_{s,\psi}^2=1/3\), \(\sigma_\psi=0\) & Tests the perfect-fluid proxy assumption & Nonzero effective anisotropic stress is required \\
Late-time expansion & Near-\(\Lambda\) residual sector with viable \(w_{\rm DE}\) & Checks whether the residual field can mimic dark energy & No viable near-\(\Lambda\) region remains \\
Internal EFT consistency & Positive post-saturation fluctuation scale and stable modes & Checks whether the threshold discharge is dynamically safe & Residual tachyonic, ghost, or gradient instability \\
\noalign{\smallskip}\hline
\end{tabular}
\end{table*}

The cleanest possible support for the PTUM proxy-fluid sector would not be a
generic excess in \(N_{\rm eff}\) alone.  It would be the combination of an
allowed extra relativistic component, the absence of the corresponding
free-streaming acoustic phase shift, and a stable effective-fluid closure with
\(\sigma_\psi=0\).  Conversely, a robust detection of an additional
neutrino-like free-streaming component would not merely tighten the pfd
abundance bound; it would directly disfavor the proxy-fluid interpretation of
that component.

%============================================================
\section{Bayesian constraints and model comparison}
\label{sec:mcmc}
%============================================================

\subsection{Likelihood setup}

We performed a Bayesian benchmark analysis using
Cobaya~\cite{Torrado2021} with the PTUM-modified CLASS backend.
The likelihood combination includes Planck~2018 high-\(\ell\)
TTTEEE-lite, low-\(\ell\) EE, CMB lensing~\cite{PlanckLike},
6dF, SDSS DR7 MGS, and SDSS DR12
BAO~\cite{Beutler2011,Ross2015,Alam2017}, and Pantheon+
supernovae~\cite{Brout2022}.

\subsection{Posterior constraints}

The marginalized PTUM constraints are
\begin{align}
\omega_b &= 0.02244^{+0.00014}_{-0.00016},\\
\omega_{\rm cdm} &= 0.1222^{+0.0013}_{-0.0023},\\
H_0 &= 68.64^{+0.55}_{-0.87}\ {\rm km\,s^{-1}\,Mpc^{-1}},\\
n_s &= 0.9639\pm0.0036,\\
\tau_{\rm reio} &= 0.0529\pm0.0071,\\
\Delta N_{\rm fluid} &< 0.31\quad(95\%\ {\rm CL}).
\label{eq:bound}
\end{align}
For the \(\Lambda\)CDM reference run:
\begin{align}
\omega_b &= 0.02235\pm0.00013,\\
\omega_{\rm cdm} &= 0.12004\pm0.00087,\\
H_0 &= 67.85\pm0.39\ {\rm km\,s^{-1}\,Mpc^{-1}},\\
n_s &= 0.9635\pm0.0037,\\
\tau_{\rm reio} &= 0.0533^{+0.0065}_{-0.0075}.
\end{align}

\begin{figure}[t]
\centering
\includegraphics[width=0.95\columnwidth]{ptum_final_triangle.pdf}
\caption{Posterior distributions from the full PTUM likelihood
analysis. The pfd amplitude peaks near zero and has a broad upper
tail, consistent with an upper-limit result.}
\label{fig:ptumtriangle}
\end{figure}

\subsection{Dataset-split analysis}

To identify which data combination drives the
\(\Delta N_{\rm fluid}\) bound, we compare results from different
dataset combinations using converged PTUM chains. The results
are summarized in Table~\ref{tab:datasplit}.

\begin{table}
\caption{Dataset-split comparison. Weighted posterior means are
shown; the last column gives the 95\% CL upper limit on
\(\Delta N_{\rm fluid}\).}
\label{tab:datasplit}
\centering
\small
\setlength{\tabcolsep}{3pt}
\begin{tabular}{@{}lcccc@{}}
\hline\noalign{\smallskip}
Data & \(H_0\) & \(\omega_c\) & \(n_s\) & \(\Delta N^{95}\) \\
\noalign{\smallskip}\hline\noalign{\smallskip}
P+B       & 68.91 & 0.1227 & 0.965 & 0.38 \\
P+L+B+S   & 68.64 & 0.1222 & 0.964 & 0.31 \\
\(\Lambda\)CDM & 67.85 & 0.1200 & 0.964 & -- \\
\noalign{\smallskip}\hline
\end{tabular}
\end{table}

Here P denotes Planck TTTEEE-lite + low-\(\ell\) EE, L denotes
Planck lensing, B denotes BAO (6dF + DR7 MGS + DR12), and S
denotes Pantheon+ supernovae.

A Planck+BAO combination already imposes a meaningful upper bound,
\(\Delta N_{\rm fluid}<0.38\) (95\% CL), while the addition of
Planck lensing and Pantheon+ tightens this to
\(\Delta N_{\rm fluid}<0.31\), a 18\% improvement. This indicates
that the PTUM pfd constraint is primarily CMB-driven, with
low-redshift probes acting mainly as geometric consistency priors
and moderate posterior sharpeners. The Hubble constant shift
is present already at the Planck+BAO level
(\(H_0=68.91\,{\rm km\,s^{-1}\,Mpc^{-1}}\)), confirming it as an
intrinsic feature of the model rather than an artifact of the
low-redshift dataset.

\subsection{Profile likelihood}

The profile likelihood of \(\Delta N_{\rm fluid}\) was constructed
from the PTUM chains by extracting the maximum posterior value in
successive \(\Delta N_{\rm fluid}\) bins. The resulting profile is
shown in Fig.\ref{fig:profile_dN}. The profile minimum lies near
\(\Delta N_{\rm fluid}\approx0.07\), indicating a mild non-zero
preference that is well within the \(1\sigma\) region. The profile
\(\Delta\chi^2=4\) level is reached near
\(\Delta N_{\rm fluid}\approx0.30\), consistent with the
marginal 95\% CL upper limit. The 68\% CL upper limit from the
marginal posterior is \(\Delta N_{\rm fluid}<0.15\).

\begin{figure}[t]
\centering
\includegraphics[width=0.95\columnwidth]{fig_profile_likelihood_dN.pdf}
\caption{Left: profile likelihood of \(\Delta N_{\rm fluid}\)
constructed from the PTUM chains. The profile \(\Delta\chi^2=4\)
level (corresponding to 95\% CL for one parameter) is reached near
\(\Delta N_{\rm fluid}\approx0.30\), consistent with the marginal
posterior bound. Right: marginal posterior PDF with 68\% and 95\%
credible regions shaded.}
\label{fig:profile_dN}
\end{figure}

\subsection{Prior robustness}

The constraint \(\Delta N_{\rm fluid}<0.31\) is data-driven rather
than prior-driven. Under flat priors Uniform\([0,0.5]\),
Uniform\([0,1]\), and Uniform\([0,3]\), the 95\% upper limit
remains at \(0.31\)--\(0.39\), unchanged at the level of the
posterior width. The dominant constraining dataset is the CMB
TT+TE+EE likelihood through Silk damping at high \(\ell\); BAO and
SN data provide geometric consistency but minimal direct sensitivity
to \(\Delta N_{\rm fluid}\).

\subsection{Model comparison}

The minimum \(\chi^2\) values are
\(\chi^2_{\rm PTUM}=2397.22\) and
\(\chi^2_{\Lambda{\rm CDM}}=2396.96\), giving
\begin{equation}
\Delta\chi^2=-0.26,\qquad \Delta{\rm AIC}=2.26.
\end{equation}
This indicates only a weak preference for the simpler model and
does not constitute a decisive exclusion of PTUM. The current data
provide a meaningful upper bound on the pfd abundance while leaving
the model observationally viable within the parameter region studied.
This interpretation is deliberately conservative: current data support
viability of the effective proxy-fluid construction, not a detection of
that sector and not a statistical preference over \(\Lambda\)CDM.

\subsection{Cosmological tensions}

PTUM shifts \(H_0\) upward by
\(+0.79\,{\rm km\,s^{-1}\,Mpc^{-1}}\) relative to \(\Lambda\)CDM
(toward the SH0ES value), reducing the Hubble tension from
\(4.7\sigma\) to \(3.5\sigma\) (25\% reduction). The \(S_8\)
tension is essentially unchanged (\(2.4\sigma\) in both models).
PTUM was not designed as a tension-solving model; these effects are
incidental consequences of the extra radiation component.

%============================================================
\section{Stability and consistency tests}
\label{sec:stability}
%============================================================

\subsection{Linear perturbation stability}

All four PTUM sectors are free of ghost and gradient instabilities.
The relevant quantities are summarized in Table~\ref{tab:stability}.

\begin{table}
\caption{Linear perturbation stability for all PTUM sectors.}
\label{tab:stability}
\centering
\begin{tabular}{lcccc}
\hline\noalign{\smallskip}
Sector & Ghost-free & \(c_s^2\) & Grad.-stable & \(\sigma\) \\
\noalign{\smallskip}\hline\noalign{\smallskip}
\(\Phi\) (creation)     & yes & 1   & yes & 0 \\
\(\psi\) (pfd)          & yes & 1/3 & yes & 0 \\
\(\chi\) (DM)           & yes & 0   & yes & 0 \\
\(\Phi_{\rm res}\) (DE) & yes & 1   & yes & 0 \\
\noalign{\smallskip}\hline
\end{tabular}
\end{table}

The pfd dispersion relation \(\omega^2=(1/3)k^2\) gives propagating,
stable acoustic oscillations for all \(k\). No tachyonic or
viscous instabilities are present in any sector.

\subsection{Isocurvature constraints}

The single-field creation mechanism produces purely adiabatic
perturbations to leading order in slow-roll. The sub-leading
isocurvature fraction from reheating asymmetry is estimated as
\(\beta_{\rm iso}\sim\epsilon^2\sim10^{-12}\)--\(10^{-6}\), which
is \(10^4\)--\(10^5\) times below the Planck 95\% CL bounds on
CDM and neutrino density isocurvature modes. The pfd isocurvature
is additionally suppressed by the small pfd fraction
(\(\Delta N/N_{\rm eff}\sim3\)--\(11\%\)).

\subsection{BBN consistency}

The PTUM bound \(\Delta N_{\rm fluid}<0.31\) is automatically
consistent with Big Bang nucleosynthesis. The BBN-only 95\%~CL
upper limit is \(\Delta N_{\rm eff}<0.40\)
(Fields et al.\ 2020~\cite{Fields2020}), which is weaker than the
CMB-derived bound. For the maximum allowed pfd abundance, the
primordial helium fraction shifts by at most
\(\delta Y_p\simeq+0.004\), corresponding to \(1.1\sigma\) of the
current \(Y_p\) measurement uncertainty.

\subsection{EFT validity}
\label{sec:eft_validity}

The PTUM EFT hierarchy \(H_*\ll V_0^{1/4}\ll M_{\rm Pl}\) is
satisfied: \(H_*/M_{\rm Pl}\sim4\times10^{-6}\) and
\(V_0/M_{\rm Pl}^4\sim10^{-10}\). All couplings are perturbative
(\(\ll4\pi\)) with margins exceeding a factor of two. For
\(\lambda\lesssim0.5\) the field excursion is super-Planckian
(a standard feature of large-field inflation models, addressed by
shift symmetry); for \(\lambda>1\) the excursion is sub-Planckian
and the EFT is fully under control without additional symmetry
input.

%============================================================
\section{Observational forecasts}
\label{sec:forecast}
%============================================================

\begin{figure*}[t]
\centering
\includegraphics[width=0.95\textwidth]{fig_forecast_PTUM.pdf}
\caption{Detection capability matrix for PTUM observables with
current and future experiments. Top left: tensor ratio \(r\) vs
\(\lambda\) with detector reach bands; top right: projected
\(\Delta N_{\rm fluid}\) upper-limit improvement; bottom left:
fluid--free-streaming phase-shift detectability; bottom right:
summary detection matrix.}
\label{fig:forecast}
\end{figure*}

The PTUM creation sector predicts a specific \(r(\lambda)\)
relation. For \(\lambda=0.02\) (\(r=0.022\)), LiteBIRD and CMB-S4
can detect the tensor signal at \(11\sigma\) and \(22\sigma\),
respectively; for \(\lambda=0.05\) (\(r=0.007\)), detection
requires CMB-S4. For \(\lambda>0.5\), \(r<10^{-3}\) is beyond
the reach of any currently planned experiment.

CMB-S4 is projected to improve the \(\Delta N_{\rm fluid}\) bound
by a factor of seven relative to the current result
(\(\sigma(\Delta N)\simeq0.027\), 95\% CL \(<0.044\)), while
CMB-HD offers a 14-fold improvement (\(<0.023\)). The fluid
versus free-streaming phase shift is detectable by CMB-S4 at
\(>3\sigma\) for \(\Delta N_{\rm fluid}\gtrsim0.20\) and at
\(>3\sigma\) by CMB-HD for \(\Delta N_{\rm fluid}\gtrsim0.10\).

Stage-IV galaxy surveys (DESI, Euclid, SKAO combined) can detect
the \(f\sigma_8\) shift at \(\sim2\)--\(3\sigma\), while full
Stage-IV combination can resolve the \(\Delta H_0=+0.79\) shift
at \(\sim5\sigma\).

\begin{center}
\textit{Appendix material.} Additional implementation notes and
consistency figures are included in the self-contained appendix below.
\end{center}

\section{Discussion}
\label{sec:discussion}
%============================================================

The main result of this work is that PTUM can now be discussed as a
quantitatively specified and observationally testable cosmological
creation framework, rather than as a purely qualitative dark-sector
scenario. The model is built around a single central idea: a
discharge-based cosmogonic event generates the hot expanding
Universe and leaves behind dark matter, dark energy, and a
relativistic proxy-fluid dark-radiation component as residual
products. In this sense, PTUM is not merely an extension of
$\Lambda$CDM by an additional radiation parameter; its primary goal
is to connect cosmogenesis and dark-sector structure within one
effective description.

Four distinct pillars have now been established:
\begin{enumerate}[label=(\roman*)]
\item the creation sector admits a plateau-type primordial potential
      compatible with current CMB constraints on $n_s$ and $r$;
\item the discharge scale is normalised to a concrete high-energy
      band, $V_0^{1/4}\sim10^{15}$--$10^{16}\,\mathrm{GeV}$;
\item the late-time dark-energy sector contains a broad viable
      region with near-$\Lambda$ behaviour;
\item the pfd sector has a closed perturbation-level
      implementation, with observational bounds established through
      Bayesian inference, profile-likelihood extraction, and
      dataset-split analysis.
\end{enumerate}

A particularly important outcome is that PTUM makes a genuinely
falsifiable prediction. The proxy-fluid component behaves as a
relativistic perfect fluid with
$w=1/3$, $c_s^2=1/3$, and $\sigma=0$, and therefore does not
produce the neutrino-like acoustic phase shift characteristic of
free-streaming dark radiation. This absence is the cleanest proposed observational discriminator
of the model. Current data constrain the abundance of this
component to $\Delta N_{\rm fluid}<0.31$ at 95\% credibility,
while future CMB-S4 and CMB-HD observations can directly test the
fluid-versus-free-streaming distinction if a non-zero dark-radiation
signal is present.

Methodologically, the numerical consistency of the framework is also
an important result. The slow-roll treatment of the creation sector
is validated by exact-background integration at the level
$\Delta n_s\sim(5.2$--$5.8)\times10^{-4}$ and
$|\Delta r|/r\sim1.3\%$--$1.5\%$, lending credibility to the
derived discharge-scale estimates and $(n_s,r)$ predictions.
Moreover, the modified \texttt{CLASS} implementation recovers
$\Lambda$CDM exactly in the limit $\Delta N_{\rm fluid}=0$, and
the perturbation sector remains ghost-free and gradient-stable
throughout the parameter region studied.

The data analysis clarifies where the present constraints originate.
The dataset-split study shows that the PTUM pfd bound is primarily
CMB-driven: the Planck$+$BAO combination alone yields
$\Delta N_{\rm fluid}<0.38$ at 95\% credibility, while adding
lensing and Pantheon$+$ tightens the limit to
$\Delta N_{\rm fluid}<0.31$. The profile likelihood is fully
consistent with the marginal posterior, with its minimum near
$\Delta N_{\rm fluid}\approx0.07$, indicating that the inferred
upper bound is not an artefact of prior volume.

At the same time, the current model-comparison result should be
interpreted with appropriate caution. PTUM is not preferred over
$\Lambda$CDM by present data:
$\Delta\chi^2=-0.26$ and $\Delta\mathrm{AIC}=2.26$ indicate a
weak preference for the simpler standard model. This outcome is not
unexpected, since PTUM introduces one additional phenomenological
degree of freedom while remaining close to $\Lambda$CDM in the
allowed region. The correct interpretation is therefore not that
PTUM is supported by current data, but that it remains observationally viable
and non-excluded within the parameter region studied while making sharper future
predictions than many dark-sector extensions.

PTUM also produces a mild upward shift in the inferred Hubble
constant,
$H_0 = 68.64^{+0.55}_{-0.87}\,\mathrm{km\,s^{-1}\,Mpc^{-1}}$,
relative to the $\Lambda$CDM fit. This reduces the SH0ES tension
from about $4.7\sigma$ to $\sim3.5\sigma$, i.e.\ by roughly
25\%, although the model was not designed as a dedicated Hubble
tension solution. By contrast, the $S_8$ tension remains
essentially unchanged. PTUM should therefore be viewed primarily as
a creation-based dark-sector framework rather than a targeted
tension-resolution model.

%============================================================
\subsection{Model summary and internal consistency}
\label{sec:summary}

PTUM's core physical picture---fuel accumulation, tachyonic
discharge, residual dark sector splitting---is internally consistent:
all sectors are ghost-free and gradient-stable (Section~8.1),
all couplings are perturbative and all energy scales satisfy
$H_*\ll M_{\rm Pl}$ (Section~8.4), and the perturbative sector
produces purely adiabatic perturbations with isocurvature fraction
$\beta_{\rm iso}\sim10^{-12}$--$10^{-6}$, far below current bounds
(Section~8.2). The BBN consistency check confirms
$\Delta N_{\rm fluid}<0.31$ is automatically compatible with the
BBN upper limit $\Delta N_{\rm eff}<0.40$ (Section~8.3).

The cosmological constant problem is re-parameterised rather than
solved: the residual dark-energy fraction
$V_{\rm res}/V_0\sim10^{-108}$--$10^{-112}$ is extremely small,
but PTUM does not provide a dynamical explanation for this hierarchy.
The magnitude remains unexplained; what PTUM provides is a concrete
interpretation of dark energy as a discharge residual and a broad parameter
region in which $w_{\rm DE}\approx-1$ is obtained.
Future data from DESI~\cite{DESI2024} and Euclid~\cite{Euclid2024}
will test whether $w_{\rm DE}$ deviates from $-1$ at percent level;
such a detection would constrain PTUM's viable region, while a
precise confirmation of $w=-1$ would leave the problem open.

Several limitations are explicitly acknowledged. The primordial
analysis uses exact background evolution rather than a full
Mukhanov--Sasaki mode-by-mode integration. The relation between
the effective discharge picture and a concrete UV completion remains
to be derived; the structural parallels with string tachyon
condensation and LQC bounce scenarios in Section~\ref{sec:uv_perspective}
should therefore be interpreted as motivation rather than derivation.
The reheating sector and the microscopic dark-matter phenomenology
of $\chi$ have not yet been analysed in comparable detail.
Finally, a nested-sampling Bayesian evidence calculation is left to
future work.

%============================================================
\subsection{Observational predictions and discriminators}
\label{sec:predictions}

The pfd sector provides the sharpest observational discriminator of
PTUM. Because $\sigma_\psi=0$, the proxy-fluid produces no
neutrino-induced CMB acoustic-peak phase shift, whereas free-streaming
dark radiation with the same energy density does. This absence is
directly testable by future CMB experiments. CMB-S4 can distinguish
fluid from free-streaming at $>3\sigma$ for
$\Delta N_{\rm fluid}\gtrsim0.20$, and CMB-HD extends this reach
to $\Delta N_{\rm fluid}\gtrsim0.10$ (Section~\ref{sec:forecast}).
The forecast matrix (Fig.\,\ref{fig:forecast}) summarises the
detection thresholds for all PTUM observables.

The creation sector predicts a specific tensor-to-scalar ratio
$r(\lambda)$. For $\lambda=0.02$ ($r\approx0.022$), LiteBIRD and
CMB-S4 can detect the signal; for $\lambda=0.05$ ($r\approx0.007$),
detection requires CMB-S4. For $\lambda>0.5$, $r<10^{-3}$ is beyond
the reach of any currently planned experiment.

On cosmological tensions: PTUM shifts $H_0$ upward by
$\sim0.8\,$km\,s$^{-1}$\,Mpc$^{-1}$, reducing the SH0ES tension
from $4.7\sigma$ to $\sim3.5\sigma$ ($\sim25\%$ improvement).
The $S_8$ tension remains essentially unchanged. This shift is an
incidental geometric effect of the extra radiation component, not a
designed tension solution.

%============================================================
\subsection{PTUM research programme: open questions and future directions}
\label{sec:roadmap}

PTUM is the first step of a broader research programme. The discharge
mechanism touches on several deep questions in fundamental physics,
each representing a natural direction for future investigation.
These directions are outlined below not as claims made in this work
but as motivated paths for the programme to pursue.

\begin{description}

\item[Singularity replacement.] 
PTUM models the hot Big Bang as a finite-energy discharge transition
rather than treating it as a primitive singular initial condition. This is analogous in spirit to
the LQC bounce~\cite{Bojowald2001,Ashtekar2011}, though PTUM operates in the classical EFT regime
($H_*/M_{\rm Pl}\sim10^{-6}$) without invoking quantum-gravity
corrections. The pre-discharge phase and its spacetime structure
remain to be specified; this requires a detailed model of the
metastable background.

\item[Dark matter phenomenology.] 
The dark-matter field $\chi$ is produced as a superheavy relic
($m_\chi\sim10^{9}$--$10^{13}\,$GeV) through parametric resonance
during the discharge. Relic abundance, detection prospects, and the
possible connection to topological defects formed at the discharge
transition are left to a dedicated study.

\item[Dark energy and the cosmological constant.] 
PTUM represents dark energy as a residual field $W(\Phi_{\rm res})$
with near-$\Lambda$ behaviour over $\sim90\%$ of the scanned
parameter space. Whether the extreme hierarchy
$V_{\rm res}/V_0\sim10^{-108}$ admits a dynamical explanation
remains an open problem; DESI and Euclid constraints on $w_0$--$w_a$
will test the viable parameter region in the near future.

\item[String theory and swampland.] 
The tachyonic activation mechanism bears structural resemblance to
open-string tachyon condensation in D-brane--anti-D-brane systems.
A full string-theoretic embedding would require specifying the
K\"ahler potential, flux background, and moduli stabilisation
procedure. Swampland criteria (distance conjecture, de Sitter
conjecture) provide additional consistency checks that PTUM
satisfies in its inflationary phase.

\item[Baryogenesis.] 
The discharge is a departure from thermal equilibrium by definition,
satisfying the third Sakharov condition. If the fuel field is extended
to include CP-violating phases, the discharge transition can generate
a net baryon number analogously to the Affleck--Dine mechanism.
Quantitative realisation requires a dedicated analysis of the
CP-violating sector.

\item[Entropy production and the arrow of time.] 
The discharge is irreversible and produces entropy. A rough estimate
gives $S_{\rm com}\sim10^{80}$ per comoving volume for the
observed cosmic entropy budget, suggesting the discharge can account
for the observed entropy without a separate low-entropy initial
condition. A detailed entropy budget and its connection to the
cosmogonically low-entropy fuel configuration remain to be
quantified.

\end{description}

In summary, PTUM reframes dark energy as a residual-field sector,
dark matter as a superheavy discharge-product candidate, and the hot
Big Bang as a finite-energy trigger event within an effective
cosmological framework. The research
programme outlined here addresses the questions that this first step
raises but does not resolve.

%============================================================
\subsection{Limitations and honest boundaries}
\label{sec:limits}

The programme must be honest about its current boundaries. PTUM
does not derive the metastable background from first principles;
does not solve the cosmological constant problem in the traditional
sense; does not include a full Mukhanov--Sasaki primordial
perturbation analysis; does not specify the reheating sector; and
does not provide a concrete UV completion. These limitations define
the scope of the present effective model.

For internal auditing, the present version of PTUM should be judged by five
separate questions rather than as a single all-or-nothing claim:
\begin{enumerate}[label=(\roman*)]
\item Does the EFT remain below its assumed cutoff and free of ghost or gradient
instabilities?
\item Does the threshold transition saturate into a locally stable post-discharge
configuration with positive curvature at the activated minimum?
\item Does the proxy-fluid sector close as a relativistic perfect fluid at the
effective linear level?
\item Do CMB, BBN, BAO, SNe, and lensing constraints leave a viable region for
$\Delta N_{\rm fluid}$ and late-time expansion parameters?
\item Can the proxy-fluid component be distinguished observationally from a
free-streaming relativistic species through the acoustic phase-shift channel?
\end{enumerate}
A failure of one test should identify a specific weak point of the model rather
than blur the status of the whole framework.  In particular, the model should
not be defended by moving a failed perturbative prediction into the background
sector or by reinterpreting a free-streaming signal as a proxy-fluid signal.

What PTUM does do is provide a controlled effective model, an implemented
Boltzmann-module prescription, a calibrated primordial scale, and a concrete
observational discriminator. The next critical step is therefore
empirical: future high-resolution CMB data will
test whether any extra radiation component behaves as a fluid
without anisotropic stress at the effective linear level or as a
free-streaming species. In this
sense, PTUM stands or falls on a clear observational criterion.

\paragraph{Falsifiability of the proxy-fluid interpretation.}
The proxy-fluid interpretation fails if future data detect an extra relativistic
component whose perturbations behave as a free-streaming species rather than as
an effective perfect fluid.  In particular, a neutrino-like acoustic phase shift
associated with the additional radiation component would disfavor the PTUM
proxy-fluid sector, whereas an extra fluid-like component with
$w_\psi=1/3$, $c_{s,\psi}^2=1/3$, and $\sigma_\psi=0$ would support its
distinctive signature.

\subsection{Final internal-consistency audit}
\label{sec:final_internal_audit}

The final version of PTUM is organized around a small set of non-interchangeable
definitions.  The audit table below states the intended meaning of each quantity
inside the present cosmological model and the interpretation that should not be
silently imported during reading or future extensions.

\begin{table*}[t]
\caption{Internal-consistency audit for the final PTUM research version.}
\label{tab:ptum_final_audit}
\centering
\begin{tabular}{p{0.17\textwidth}p{0.36\textwidth}p{0.35\textwidth}}
\hline\noalign{\smallskip}
Item & PTUM meaning & Guardrail \\
\noalign{\smallskip}\hline\noalign{\smallskip}
$\Phi$ & Quantum fuel field controlling the discharge threshold & An effective scalar variable, not a completed microscopic derivation of the pre-discharge background \\
$\psi$ & Activated proxy mode whose perturbations are represented as fluid-like dark radiation & Not a free-streaming neutrino-like hierarchy unless $\sigma_\psi\neq0$ is explicitly introduced \\
$\chi$ & Passive post-discharge dark-matter candidate sector & Not used as evidence for a specific particle-identification claim \\
$\Phi_{\rm res}$ & Residual near-$\Lambda$ dark-energy-like component & Not a solution of the cosmological-constant problem by itself \\
$\Delta N_{\rm fluid}$ & Effective abundance of fluid-like dark radiation & Not identical to arbitrary $\Delta N_{\rm eff}$ unless the perturbative closure is specified \\
Operational closure & Tachyonic activation, local post-discharge stability, and perfect-fluid perturbative closure & Not a UV completion or a proof that this is the unique possible potential form \\
\noalign{\smallskip}\hline
\end{tabular}
\end{table*}

The final claim boundary is therefore narrow but useful: PTUM is a testable
effective phase-transition cosmology model whose sharpest discriminator is a
fluid-like relativistic component with no free-streaming acoustic phase shift.
It is not, at the present stage, a UV-complete cosmogony, a derivation of the
metastable pre-discharge background, or a solution to the cosmological-constant
hierarchy.  This restricted formulation is the version that should be carried
forward as the PTUM library-final manuscript.

%============================================================
\section{Conclusion}
\label{sec:conclusion}
%============================================================

We have presented the Phase Transition Universes Model as a
discharge-based cosmological creation framework in which a single
effective quantum-fuel discharge provides a phenomenological description
of the onset of a hot expanding cosmological phase and sources a
dark-matter candidate sector, a residual dark-energy sector, and a
proxy-fluid dark-radiation sector as residual components. The analysis was carried from the
level of effective field theory and background dynamics to
perturbations, Boltzmann-code implementation, and benchmark comparison
with current cosmological data.

The main conclusions of this work are as follows:

\begin{enumerate}[label=(\arabic*)]
\item The PTUM fuel potential is a plateau-type creation potential
      capable of supporting a primordial accelerated phase
      consistent with present constraints on $(n_s,r)$.

\item Exact-background calculations validate the slow-roll treatment
      across $\lambda=0.02$--$5$, with
      $\Delta n_s\sim(5.2$--$5.8)\times10^{-4}$ and
      $|\Delta r|/r\sim1.3\%$--$1.5\%$.

\item Matching the observed scalar amplitude fixes the creation
      scale to
      $V_0^{1/4}\sim 2\times10^{15}$--$1.3\times10^{16}\,
      \mathrm{GeV}$, corresponding to
      $H_*\sim10^{12}$--$10^{13}\,\mathrm{GeV}$.

\item The late-time residual potential admits a broad viable region,
      with $1+w_{\rm DE}<0.01$ over 89.6\% of the scanned
      $(\kappa,\lambda_{\rm res})$ space, so that PTUM can
      reproduce $w_{\rm DE}\approx -1$ over most of its allowed
      parameter domain.

\item The proxy-fluid dark-radiation component has the distinctive
      effective-linear perturbative signature
      $w=1/3$, $c_s^2=1/3$, and $\sigma=0$. Unlike free-streaming
      dark radiation, it therefore produces no neutrino-induced CMB
      acoustic phase shift. This is the principal observational
      discriminator of PTUM.

\item A Bayesian benchmark analysis using Planck~2018 TT,TE,EE-lite,
      low-$\ell$ EE, CMB lensing, BAO, and Pantheon$+$ yields
      $\Delta N_{\rm fluid}<0.31$ at 95\% credibility. This bound
      is automatically consistent with BBN constraints.

\item The dataset-split analysis shows that the pfd bound is mainly
      CMB-driven: Planck$+$BAO alone gives
      $\Delta N_{\rm fluid}<0.38$, while the full combined data
      set tightens the limit to $\Delta N_{\rm fluid}<0.31$.
      The profile likelihood is consistent with the marginal
      posterior and reaches the $\Delta\chi^2=4$ level near
      $\Delta N_{\rm fluid}\approx0.30$.

\item Relative to $\Lambda$CDM, PTUM gives
      $\Delta\chi^2=-0.26$ and
      $\Delta\mathrm{AIC}=2.26$. Present data therefore show a
      weak preference for $\Lambda$CDM, but do not exclude PTUM.
      The model remains observationally viable.

\item All sectors are perturbatively stable and EFT-consistent:
      the theory is ghost-free, gradient-stable, adiabatic at
      leading order, compatible with isocurvature bounds, and valid
      in the regime $H_* \ll M_{\rm Pl}$ with perturbative
      couplings.

\item Future observations provide a forecast-level test. CMB-S4 can
      improve the $\Delta N_{\rm fluid}$ sensitivity by roughly a
      factor of $7$ and may distinguish fluid from free-streaming dark
      radiation at high significance if
      $\Delta N_{\rm fluid}\gtrsim0.20$, while CMB-HD can extend
      this reach to $\Delta N_{\rm fluid}\gtrsim0.10$.
\item The discharge-sector closure reduces to a compact EFT condition:
      the tachyonic threshold activates the proxy phase, the activated minimum
      is locally stable, and the resulting coherent component is tested through
      the effective-linear proxy-fluid signature $w_\psi=1/3$,
      $c_{s,\psi}^2=1/3$, and $\sigma_\psi=0$.
\end{enumerate}

Taken together, these results support PTUM as a viable effective
phase-transition cosmology framework with a falsifiable proxy-fluid
dark-radiation channel. Its present value lies not in outperforming
$\Lambda$CDM with current data, but in organizing a discharge-based dark-sector
scenario around a concrete perturbative discriminator. If future CMB
observations reveal extra relativistic energy density without the corresponding
free-streaming phase shift, PTUM-like fluid dark radiation would become a
natural effective interpretation. If they do not, the allowed proxy-fluid
parameter space will be sharply constrained. In either case, the framework
remains in a well-defined predictive regime.

\begin{acknowledgement}
The author acknowledges the developers of
CLASS~\cite{Blas2011,Lesgourgues2011} and
Cobaya~\cite{Torrado2021} for providing open-source cosmological
analysis tools. During the preparation of this work the author used
large language model tools to assist with source-code debugging,
Python scripting, and \LaTeX{} formatting. After using these tools,
the author reviewed and edited the content as needed and takes full
responsibility for the content of the publication.
\end{acknowledgement}

%============================================================
% APPENDIX
%============================================================

\appendix

%------------------------------------------------------------
\appendix

\section{CLASS implementation details}
\label{app:class}
%------------------------------------------------------------

The observational predictions of PTUM are computed using a modified
version of the Cosmic Linear Anisotropy Solving System
(\texttt{CLASS}~v3.2.0)~\cite{Blas2011}. This appendix documents
the implementation to ensure reproducibility.

\subsection{Background module}

The proxy-fluid dark radiation component is implemented as an
additional perfect fluid with equation of state $w_\psi=1/3$ and
energy density parametrised by the effective number of fluid
degrees of freedom $\Delta N_{\rm fluid}$. The energy density at
the reference epoch is
\begin{equation}
\rho_\psi
= \frac{7}{8}\left(\frac{4}{11}\right)^{4/3}
  \Delta N_{\rm fluid}\,\rho_\gamma\,,
\label{eq:rho_pfd_app}
\end{equation}
where $\rho_\gamma$ is the photon energy density. This component
scales as $\rho_\psi\propto a^{-4}$ and is added to the Friedmann
equation alongside standard radiation, matter, and dark energy.

\subsection{Perturbation module}

At the perturbation level the pfd component is treated as a
\emph{perfect fluid} with
\begin{equation}
c_s^2 = \tfrac{1}{3}\,,\qquad \sigma_\psi = 0\,,
\label{eq:pfd_pert_app}
\end{equation}
where $c_s^2$ is the rest-frame sound speed and $\sigma_\psi$ is
the anisotropic stress. This is the key distinction from
free-streaming dark radiation (e.g.\ additional neutrino species),
which carries non-zero anisotropic stress. The perturbation
equations in synchronous gauge are
\begin{align}
\dot{\delta}_\psi
  &= -\tfrac{4}{3}\bigl(\theta_\psi+\tfrac{\dot{h}}{2}\bigr)\,,
\label{eq:delta_pfd_app}\\[4pt]
\dot{\theta}_\psi
  &= \tfrac{k^2}{4}\,\delta_\psi\,,
\label{eq:theta_pfd_app}
\end{align}
corresponding to a relativistic perfect fluid with no higher
multipoles.

\subsection{State-vector layout}

The pfd component is assigned a dedicated index in the
\texttt{CLASS} state vector, positioned after the standard
neutrino hierarchy. Initial conditions are set in the
radiation-dominated era assuming adiabatic perturbations:
\begin{equation}
\delta_\psi^{(\rm ini)} = \delta_\gamma^{(\rm ini)}\,,
\qquad
\theta_\psi^{(\rm ini)} = \theta_\gamma^{(\rm ini)}\,.
\end{equation}
No additional free-streaming hierarchy is introduced for the proxy-fluid
component. In the implementation, the pfd sector is evolved with density and
velocity perturbations only, while its higher multipoles are absent by
construction. This is not a numerical shortcut: it is the effective linear
closure that defines the PTUM proxy-fluid discriminator.

\subsection{Validation}

The implementation has been validated through the following checks:

\begin{enumerate}
\item \textbf{$\Lambda$CDM recovery:} Setting
      $\Delta N_{\rm fluid}=0$ recovers the standard \texttt{CLASS}
      $\Lambda$CDM output to machine precision in both
      $C_\ell^{TT}$ and $P(k)$.

\item \textbf{Background consistency:} The Friedmann constraint
      $\sum_i\Omega_i=1$ is satisfied to $<10^{-12}$ at all
      redshifts.

\item \textbf{Perturbation stability:} No numerical instabilities
      arise for $\Delta N_{\rm fluid}\in[0,1]$ across the full
      range of angular scales $\ell\in[2,5000]$.

\item \textbf{Comparison with \texttt{ncdm} module:} The pfd
      spectrum differs from the free-streaming
      $\Delta N_{\rm eff}$ spectrum in the expected manner:
      identical background evolution but suppressed phase shift
      in $C_\ell^{TT}$ due to $\sigma=0$.
\end{enumerate}

\subsection{Code availability}
\label{app:code-availability}

The manuscript source, final figure artifacts, claim-boundary notes,
and reproducibility scaffold are publicly available at
\url{https://github.com/ErdemEsa/ptum-cosmology.git}. The repository
includes the public PTUM code-release structure and a clean
\texttt{CLASS}~v3.2.0 patch for the proxy-fluid dark-radiation sector.
The patch introduces the effective-fluid parameter
$\Delta N_{\rm fluid}$, the associated background contribution, and the
linear perturbation variables for a perfect-fluid dark-radiation
component with $w=1/3$, $c_s^2=1/3$, and $\sigma=0$. It was generated as
a path-limited diff against a clean \texttt{CLASS}~v3.2.0 baseline and
has been checked with \texttt{git apply --check} against the corresponding
clean source tree. Thus, the public repository now provides the
manuscript archive and the PTUM-modified \texttt{CLASS} patch needed to
inspect and re-apply the core Boltzmann-code implementation.

The production-level \texttt{Cobaya} configuration files and full MCMC
chains are not included in this manuscript ZIP. They will be added to the
same public repository or archived externally in a later reproducibility
release. Until those production configuration files and chains are
released, the Bayesian constraints reported here should be read as
benchmark-level summaries rather than as a fully public end-to-end
pipeline release.

%------------------------------------------------------------
\section{Appendix figures and consistency diagnostics}
\label{app:figures}
%------------------------------------------------------------

The following figures are included in this appendix because they
directly support the main-text consistency analysis. The library-final
version is therefore self-contained and does not require the obsolete
external supplementary-material pointer.

%--- Appendix Figure A1 ---
\begin{figure}[t]
\centering
\includegraphics[width=\columnwidth]{fig_BBN_consistency.pdf}
\caption{BBN consistency check for PTUM. The shaded region shows
the BBN-allowed range $\Delta N_{\rm eff}<0.40$ from primordial
helium and deuterium abundances~\cite{Fields2020,Cyburt2016}. The
vertical line marks the PTUM 95\% upper limit
$\Delta N_{\rm fluid}<0.31$, which lies well within the
BBN-allowed region. At the maximum allowed pfd abundance, the
predicted shift in $Y_p$ is approximately $1.1\sigma$.}
\label{fig:bbn}
\end{figure}

%--- Appendix Figure A2 ---
\begin{figure*}[t]
\centering
\includegraphics[width=\textwidth]{fig_fluid_vs_freestream.pdf}
\caption{Comparison of CMB $TT$ (left) and $EE$ (right) power
spectrum residuals for PTUM proxy-fluid dark radiation
($\sigma=0$, solid) and free-streaming dark radiation
($\sigma\neq0$, dashed) at $\Delta N=0.10$. The free-streaming
component produces a characteristic phase shift of the acoustic
peaks, while the proxy-fluid shifts peak amplitudes without
displacing their positions. The distinction is especially
pronounced in $EE$ at $\ell>1000$.}
\label{fig:fluid_vs_fs}
\end{figure*}

%--- Appendix Figure A3 ---
\begin{figure*}[t]
\centering
\includegraphics[width=\textwidth]{lcdm_vs_ptum_triangle.pdf}
\caption{Comparison of marginalised posterior distributions for
$\Lambda$CDM (blue) and PTUM (red) from the full data
combination. The standard parameters are consistent between the
two models. PTUM shows a mild broadening in $\omega_{\rm cdm}$
and $H_0$ due to the degeneracy with $\Delta N_{\rm fluid}$, and
a systematic upward shift of $H_0$ by
$\sim0.8\;\mathrm{km\,s^{-1}\,Mpc^{-1}}$.}
\label{fig:lcdm_vs_ptum}
\end{figure*}

%--- Appendix Figure A4 ---
\begin{figure}[t]
\centering
\includegraphics[width=\columnwidth]{fig_perturbation_stability.pdf}
\caption{Linear perturbation stability diagnostics for all PTUM
sectors. Top: squared sound speed $c_s^2$ as a function of scale
factor, confirming $c_s^2>0$ (no gradient instability) throughout.
Bottom: effective squared mass confirming the absence of tachyonic
instabilities. All sectors are ghost-free by construction.}
\label{fig:stability}
\end{figure}

%--- Appendix Figure A5 ---
\begin{figure}[t]
\centering
\includegraphics[width=\columnwidth]{fig_eft_validity.pdf}
\caption{EFT validity diagnostics for PTUM. The inflationary
Hubble scale $H_*$, potential energy $V_0^{1/4}$, slow-roll
parameters $\epsilon$ and $|\eta|$, and coupling constants are
shown as functions of $\lambda$. All quantities satisfy EFT
consistency: $H_*\ll M_{\rm Pl}$, $V_0<M_{\rm Pl}^4$,
$\epsilon,|\eta|\ll1$, and perturbative couplings. For
$\lambda\lesssim0.05$ the field excursion is super-Planckian,
protectable by an approximate shift
symmetry~\cite{Freese1990,Silverstein2008}.}
\label{fig:eft}
\end{figure}

%============================================================
% REFERENCES
%============================================================

\begin{thebibliography}{99}

% ---- Cosmological parameters / Planck ----
\bibitem{Planck2018params}
N.~Aghanim et al.\ [Planck Collaboration],
Astron.\ Astrophys.\ \textbf{641}, A6 (2020).

% ---- Quintessence / dynamical dark energy ----
\bibitem{RatraPeebles}
B.~Ratra, P.~J.~E.~Peebles,
Phys.\ Rev.\ D \textbf{37}, 3406 (1988).

\bibitem{Caldwell1998}
R.~R.~Caldwell, R.~Dave, P.~J.~Steinhardt,
Phys.\ Rev.\ Lett.\ \textbf{80}, 1582 (1998).

% ---- Unified dark sector ----
\bibitem{Kamenshchik2001}
A.~Kamenshchik, U.~Moschella, V.~Pasquier,
Phys.\ Lett.\ B \textbf{511}, 265 (2001).

\bibitem{Bento2002}
M.~C.~Bento, O.~Bertolami, A.~A.~Sen,
Phys.\ Rev.\ D \textbf{66}, 043507 (2002).

% ---- Neff / BBN ----
\bibitem{Mangano2005}
G.~Mangano, G.~Miele, S.~Pastor, T.~Pinto, O.~Pisanti,
P.~D.~Serpico,
Nucl.\ Phys.\ B \textbf{729}, 221 (2005).

\bibitem{Cyburt2016}
R.~H.~Cyburt, B.~D.~Fields, K.~A.~Olive, T.-H.~Yeh,
Rev.\ Mod.\ Phys.\ \textbf{88}, 015004 (2016).

% ---- CMB phase shift / dark radiation ----
\bibitem{Bashinsky2004}
S.~Bashinsky, U.~Seljak,
Phys.\ Rev.\ D \textbf{69}, 083002 (2004).

\bibitem{Baumann2016}
D.~Baumann, D.~Green, J.~Meyers, B.~Wallisch,
JCAP \textbf{01}, 007 (2016).

\bibitem{Follin2015}
B.~Follin, L.~Knox, M.~Millea, Z.~Pan,
Phys.\ Rev.\ Lett.\ \textbf{115}, 091301 (2015).

\bibitem{CyrRacine2014}
F.-Y.~Cyr-Racine, K.~Sigurdson,
Phys.\ Rev.\ D \textbf{87}, 103515 (2013).

% ---- Boltzmann / MCMC tools ----
\bibitem{Blas2011}
D.~Blas, J.~Lesgourgues, T.~Tram,
JCAP \textbf{07}, 034 (2011).

\bibitem{Lesgourgues2011}
J.~Lesgourgues,
arXiv:1104.2932.

\bibitem{Torrado2021}
J.~Torrado, A.~Lewis,
JCAP \textbf{05}, 057 (2021).

% ---- Perturbation theory / field theory ----
\bibitem{MaBertschinger}
C.-P.~Ma, E.~Bertschinger,
Astrophys.\ J.\ \textbf{455}, 7 (1995).

\bibitem{Turner1983}
M.~S.~Turner,
Phys.\ Rev.\ D \textbf{28}, 1243 (1983).

\bibitem{ColemanWeinberg}
S.~R.~Coleman, E.~J.~Weinberg,
Phys.\ Rev.\ D \textbf{7}, 1888 (1973).

\bibitem{Maldacena2003}
J.~M.~Maldacena,
JHEP \textbf{05}, 013 (2003).

% ---- Observational data ----
\bibitem{PlanckLike}
N.~Aghanim et al.\ [Planck Collaboration],
Astron.\ Astrophys.\ \textbf{641}, A5 (2020).

\bibitem{PlanckNG}
Y.~Akrami et al.\ [Planck Collaboration],
Astron.\ Astrophys.\ \textbf{641}, A9 (2020).

\bibitem{Beutler2011}
F.~Beutler et al.,
Mon.\ Not.\ R.\ Astron.\ Soc.\ \textbf{416}, 3017 (2011).

\bibitem{Ross2015}
A.~J.~Ross et al.,
Mon.\ Not.\ R.\ Astron.\ Soc.\ \textbf{449}, 835 (2015).

\bibitem{Alam2017}
S.~Alam et al.,
Mon.\ Not.\ R.\ Astron.\ Soc.\ \textbf{470}, 2617 (2017).

\bibitem{Brout2022}
D.~Brout et al.,
Astrophys.\ J.\ \textbf{938}, 110 (2022).

\bibitem{BICEPKeck2021}
P.~A.~R.~Ade et al.\ [BICEP/Keck Collaboration],
Phys.\ Rev.\ Lett.\ \textbf{127}, 151301 (2021).

\bibitem{Fields2020}
B.~D.~Fields, K.~A.~Olive, T.-H.~Yeh, C.~Young,
JCAP \textbf{03}, 010 (2020).

% ---- Experimental programmes ----
\bibitem{CMBS4}
K.~Abazajian et al.,
arXiv:1610.02743.

\bibitem{DESI2024}
DESI Collaboration,
arXiv:2404.03000.

\bibitem{Euclid2024}
Euclid Collaboration,
Astron.\ Astrophys.\ \textbf{642}, A191 (2020).

% ---- Cosmological frameworks ----
\bibitem{Susskind2003}
L.~Susskind,
arXiv:hep-th/0302219.

\bibitem{Smolin1997}
L.~Smolin,
\textit{The Life of the Cosmos}
(Oxford University Press, Oxford, 1997).

% ---- UV perspective: string theory ----
\bibitem{Sen2002}
A.~Sen,
JHEP \textbf{0207}, 065 (2002).

\bibitem{Dvali1999}
G.~Dvali, S.-H.~H.~Tye,
Phys.\ Lett.\ B \textbf{450}, 72 (1999).

\bibitem{KKLMMT2003}
S.~Kachru, R.~Kallosh, A.~Linde, J.~Maldacena, L.~McAllister,
S.~P.~Trivedi,
JCAP \textbf{0310}, 013 (2003).

% ---- UV perspective: swampland ----
\bibitem{Obied2018}
G.~Obied, H.~Ooguri, L.~Spodyneiko, C.~Vafa,
arXiv:1806.08362.

\bibitem{Ooguri2019}
H.~Ooguri, E.~Palti, G.~Shiu, C.~Vafa,
Phys.\ Lett.\ B \textbf{788}, 180 (2019).

% ---- UV perspective: shift symmetry ----
\bibitem{Freese1990}
K.~Freese, J.~A.~Frieman, A.~V.~Olinto,
Phys.\ Rev.\ Lett.\ \textbf{65}, 3233 (1990).

\bibitem{Silverstein2008}
E.~Silverstein, A.~Westphal,
Phys.\ Rev.\ D \textbf{78}, 106003 (2008).

% ---- LQC Bounce ----

\bibitem{Bojowald2001}
M.~Bojowald,
Phys.\ Rev.\ Lett.\ \textbf{86}, 5227 (2001).

\bibitem{Ashtekar2011}
A.~Ashtekar, P.~Singh,
Class.\ Quantum Grav.\ \textbf{28}, 213001 (2011).

% ---- Black hole information paradox ----

\bibitem{Hawking1975}
S.~W.~Hawking,
Commun.\ Math.\ Phys.\ \textbf{43}, 199 (1975).

\bibitem{Hawking1976}
S.~W.~Hawking,
Phys.\ Rev.\ D \textbf{14}, 2460 (1976).

% ---- LZ direct detection ----

\bibitem{LZ2022}
LZ Collaboration,
Phys.\ Rev.\ Lett.\ \textbf{131}, 041002 (2023).

\end{thebibliography}

\end{document}